
\begin{figure}[!htbp]
\centering
\begin{subfigure}[t]{0.49\textwidth}\centering
\includegraphics[width=\linewidth]{figures/ch4/cp112_generated/pp_base_shear.png}\subcaption{Base shear}
\end{subfigure}\hfill
\begin{subfigure}[t]{0.49\textwidth}\centering
\includegraphics[width=\linewidth]{figures/ch4/cp112_generated/pp_displacement.png}\subcaption{Roof displacement}
\end{subfigure}
\caption{Performance-point response comparison.}
\label{fig:pp_model_comparison_cp112}
\end{figure}

Figure~\ref{fig:pp_model_comparison_cp112} shows that the model governing performance-point base shear is not the same as the model governing performance-point roof displacement. The Regular model has the largest performance-point base shear, whereas the U-shaped model has the largest roof displacement, reinforcing the need to interpret force and deformation demand separately.

\FloatBarrier

\subsection{Hinge-Evidence Summary}
The detailed final-step hinge screenshots remain in Appendix~\ref{appx:hinge_formation_results}. All seven models have X- and Y-direction final-step images for the five matched demand cases. However, the available analysis records do not provide a member-level table identifying a unique ``critical hinge location'', beam/column classification, and final hinge state for every model/direction. Table~\ref{tab:hinge_evidence_cp112} therefore reports the evidence status rather than inventing those fields from screenshots.
\begin{table}[!htbp]
\centering
\caption{Hinge evidence available for the seven-model comparison.}
\label{tab:hinge_evidence_cp112}
\small
\resizebox{\textwidth}{!}{%
\begin{tabular}{@{}lccp{6.3cm}@{}}
\toprule
Model & X final-step image & Y final-step image & Member-level result status \\
\midrule
Regular & Available & Available & Critical member/location/state not available in tabulated results \\
L & Available & Available & Same limitation \\
L1 & Available & Available & Same limitation \\
L2 & Available & Available & Same limitation \\
T & Available & Available & Same limitation \\
U & Available & Available & Same limitation \\
Non-Parallel & Available & Available & Same limitation \\
\bottomrule
\end{tabular}
}
\end{table}
